\textbf{Phần 1. Trắc nghiệm lựa chọn}
\setcounterpageref{ex}{0}

\begin{ex} %[2X13N65]
Biểu đồ tranh dùng cái gì để biểu diễn số liệu?
\choice
{Các đoạn thẳng nối tiếp}
{Các hình khối chữ nhật}
{\True Các hình ảnh hoặc biểu tượng tranh vẽ}
{Các con số đơn thuần}
\loigiai{
Biểu đồ tranh dùng các hình ảnh hoặc biểu tượng.\\
\textbf{Đáp án của câu hỏi là C}
}
\end{ex}

\begin{ex} %[2X13N65]
Nếu mỗi biểu tượng ngôi sao $\bigstar$ biểu diễn cho 1 bạn học sinh, có 5 ngôi sao $\bigstar$ tương ứng với
\choice
{$1$ bạn}
{\True $5$ bạn học sinh}
{$10$ bạn}
{$15$ bạn}
\loigiai{
5 ngôi sao tương ứng với 5 bạn học sinh.\\
\textbf{Đáp án của câu hỏi là B}
}
\end{ex}

\begin{ex} %[2X13H65]
Nếu mỗi biểu tượng quả bóng biểu diễn cho 2 quả bóng thực tế, có 4 biểu tượng tương ứng với
\choice
{$6$ quả bóng}
{$10$ quả bóng}
{\True $8$ quả bóng}
{$4$ quả bóng}
\loigiai{
$2 \times 4 = 8$ quả bóng.\\
\textbf{Đáp án của câu hỏi là C}
}
\end{ex}

\begin{ex} %[2X13H65]
Nếu mỗi bông hoa biểu diễn cho 5 bông hoa thật, có 3 bông hoa biểu tượng tương ứng với
\choice
{\True $15$ bông hoa}
{$10$ bông}
{$20$ bông}
{$25$ bông}
\loigiai{
$5 \times 3 = 15$ bông hoa.\\
\textbf{Đáp án của câu hỏi là A}
}
\end{ex}

\begin{ex} %[2X13H65]
Quan sát biểu đồ tranh về số quả táo của 3 bạn: Mai có 4 quả táo, Lan có 6 quả táo, Hà có 3 quả táo (mỗi hình vẽ là 1 quả). Bạn có nhiều táo nhất là
\choice
{Mai}
{Hà}
{Các bạn bằng nhau}
{\True Lan ($6$ quả)}
\loigiai{
Lan có 6 quả táo (nhiều nhất).\\
\textbf{Đáp án của câu hỏi là D}
}
\end{ex}

\begin{ex} %[2X13H65]
Lan có nhiều hơn Hà số quả táo là
\choice
{\True $3$ quả}
{$2$ quả}
{$4$ quả}
{$5$ quả}
\loigiai{
$6 - 3 = 3$ quả.\\
\textbf{Đáp án của câu hỏi là A}
}
\end{ex}

\begin{ex} %[2X13V65]
Tổng số quả táo của cả ba bạn Mai, Lan, Hà là
\choice
{$12$ quả}
{\True $13$ quả}
{$14$ quả}
{$15$ quả}
\loigiai{
$4 + 6 + 3 = 13$ quả táo.\\
\textbf{Đáp án của câu hỏi là B}
}
\end{ex}

\begin{ex} %[2X13V65]
Biểu đồ tranh số cây trồng của các tổ (mỗi hình cây $= 2$ cây thật): Tổ 1 có 3 hình cây, Tổ 2 có 5 hình cây. Tổ 2 trồng được số cây thật là
\choice
{$8$ cây}
{$12$ cây}
{$15$ cây}
{\True $10$ cây}
\loigiai{
$2 \times 5 = 10$ cây.\\
\textbf{Đáp án của câu hỏi là D}
}
\end{ex}

\begin{ex} %[2X13V65]
Tổ 2 trồng nhiều hơn Tổ 1 số cây thật là
\choice
{\True $4$ cây}
{$2$ cây}
{$6$ cây}
{$8$ cây}
\loigiai{
Tổ 1: $2 \times 3 = 6$ cây. Tổ 2: 10 cây.\\
$10 - 6 = 4$ cây.\\
\textbf{Đáp án của câu hỏi là A}
}
\end{ex}

\begin{ex} %[2X13V65]
Biểu đồ tranh số thú nuôi trong trang trại: Bò có 4 con biểu tượng, Cừu có 6 con biểu tượng, Lợn có 5 con biểu tượng (mỗi biểu tượng $= 5$ con thật). Số cừu trong trang trại là
\choice
{$25$ con}
{$35$ con}
{$40$ con}
{\True $30$ con}
\loigiai{
$5 \times 6 = 30$ con cừu.\\
\textbf{Đáp án của câu hỏi là D}
}
\end{ex}

\begin{ex} %[2X13V65]
Tổng số con vật có trong trang trại đó là
\choice
{$70$ con}
{$80$ con}
{\True $75$ con}
{$85$ con}
\loigiai{
Tổng số biểu tượng: $4 + 6 + 5 = 15$.\\
Tổng số con: $5 \times 15 = 75$ con.\\
\textbf{Đáp án của câu hỏi là C}
}
\end{ex}

\begin{ex} %[2X13C65]
Trên biểu đồ tranh, bạn An có 4 hình trái tim, bạn Bình có 7 hình trái tim. Mỗi hình trái tim biểu diễn cho 2 cuốn sách. Hỏi Bình có nhiều hơn An bao nhiêu cuốn sách?
\choice
{$4$ cuốn}
{\True $6$ cuốn}
{$8$ cuốn}
{$10$ cuốn}
\loigiai{
Bình hơn An: $7 - 4 = 3$ hình trái tim.\\
Số sách nhiều hơn: $2 \times 3 = 6$ cuốn.\\
\textbf{Đáp án của câu hỏi là B}
}
\end{ex}

\textbf{Phần 2. Trắc nghiệm đúng sai}
\setcounterpageref{ex}{0}

\begin{ex} %[2X13H65]
Biểu đồ tranh ghi số ô tô bán được trong tuần (mỗi ô tô vẽ $= 2$ xe thật): Thứ Hai có 3 hình vẽ, Thứ Ba có 4 hình vẽ. Xét tính đúng sai:
\choiceTF
{\True Thứ Hai bán được 6 xe ô tô thật ($2 \times 3 = 6$).}
{\True Thứ Ba bán được 8 xe ô tô thật ($2 \times 4 = 8$).}
{Thứ Ba bán nhiều hơn Thứ Hai 1 xe thật.}
{\True Thứ Ba bán nhiều hơn Thứ Hai 2 xe thật ($8 - 6 = 2$).}
\loigiai{
\begin{itemize}
    \item Thứ Hai: 6 xe: Đúng.
    \item Thứ Ba: 8 xe: Đúng.
    \item Hơn 2 xe (không phải 1 xe). Ý 3 Sai, Ý 4 Đúng.
\end{itemize}
\begin{itemize}
    \item a) \textbf{Đúng}.
    \item b) \textbf{Đúng}.
    \item c) \textbf{Sai}.
    \item d) \textbf{Đúng}.
\end{itemize}
}
\end{ex}

\begin{ex} %[2X13V65]
Biểu đồ tranh về số quả bóng (mỗi biểu tượng $= 5$ quả): Lớp 2A có 4 biểu tượng, Lớp 2B có 3 biểu tượng. Xét tính đúng sai:
\choiceTF
{\True Lớp 2A có 20 quả bóng ($5 \times 4 = 20$).}
{\True Lớp 2B có 15 quả bóng ($5 \times 3 = 15$).}
{Cả hai lớp có tất cả 40 quả bóng.}
{\True Cả hai lớp có tất cả 35 quả bóng ($20 + 15 = 35$).}
\loigiai{
\begin{itemize}
    \item Lớp 2A có 20 quả: Đúng.
    \item Lớp 2B có 15 quả: Đúng.
    \item Tổng: $20 + 15 = 35 \neq 40$. Ý 3 Sai, Ý 4 Đúng.
\end{itemize}
\begin{itemize}
    \item a) \textbf{Đúng}.
    \item b) \textbf{Đúng}.
    \item c) \textbf{Sai}.
    \item d) \textbf{Đúng}.
\end{itemize}
}
\end{ex}

\textbf{Phần 3. Trả lời ngắn}
\setcounterpageref{ex}{0}

\begin{ex} %[2X13H65]
Trên biểu đồ tranh, mỗi ngôi sao biểu diễn cho 2 điểm mười. Bạn Nam có 5 ngôi sao. Hỏi Nam có bao nhiêu điểm mười?
\shortans{10}
\loigiai{
\textbf{Đáp số:} $10$\\
$2 \times 5 = 10$ điểm mười.
}
\end{ex}

\begin{ex} %[2X13H65]
Bạn Mai có 7 ngôi sao (mỗi ngôi sao $= 2$ điểm mười). Hỏi Mai có bao nhiêu điểm mười?
\shortans{14}
\loigiai{
\textbf{Đáp số:} $14$\\
$2 \times 7 = 14$ điểm mười.
}
\end{ex}

\begin{ex} %[2X13V65]
Trên biểu đồ, mỗi hình cái cây biểu diễn cho 5 cây thật. Lớp 2A có 6 hình cái cây. Hỏi lớp 2A trồng được bao nhiêu cây?
\shortans{30}
\loigiai{
\textbf{Đáp số:} $30$\\
$5 \times 6 = 30$ cây.
}
\end{ex}

\begin{ex} %[2X13V65]
Lớp 2B có 8 hình cái cây (mỗi hình $= 5$ cây). Hỏi lớp 2B trồng được bao nhiêu cây?
\shortans{40}
\loigiai{
\textbf{Đáp số:} $40$\\
$5 \times 8 = 40$ cây.
}
\end{ex}

\begin{ex} %[2X13V65]
Lớp 2B trồng nhiều hơn lớp 2A bao nhiêu cây?
\shortans{10}
\loigiai{
\textbf{Đáp số:} $10$\\
$40 - 30 = 10$ cây.
}
\end{ex}

\begin{ex} %[2X13C65]
Có 45 quyển sách cần biểu diễn lên biểu đồ tranh. Biết mỗi biểu tượng quyển sách biểu diễn cho 5 cuốn sách thật. Hỏi cần vẽ bao nhiêu biểu tượng?
\shortans{9}
\loigiai{
\textbf{Đáp số:} $9$\\
$45 : 5 = 9$ biểu tượng.
}
\end{ex}
